\chapter{Introduction générale}

\section{Point de départ}

Une imprimante industrielle peut paraître secondaire à côté d'un four, d'un
robot de chargement ou d'une ligne d'assemblage. Pourtant, si elle ne marque
plus correctement le verre, la pièce perd une partie de sa traçabilité et la
production peut être arrêtée. Dans le cas d'une imprimante à jet d'encre
continu, les causes sont variées : niveau de solvant trop bas, viscosité hors
plage, défaut de pression, buse obstruée ou consommable périmé.

Sur la ligne étudiée, les automates communiquent en PROFINET et les données
remontent déjà vers Litmus Edge. Elles restent toutefois difficiles à
exploiter par les personnes qui en ont besoin. L'opérateur veut connaître
l'état immédiat de la machine. La maintenance veut comprendre les dérives et
les défauts récurrents. La direction regarde surtout la disponibilité et le
temps perdu. Ces trois besoins ne se traduisent pas par le même écran.

Le chapitre consacré à l'état de l'art situe cette difficulté parmi les
approches de supervision industrielle avant de formuler la problématique
retenue.

\section{Objectifs}

Le projet poursuit les objectifs suivants :

\begin{itemize}
  \item afficher l'état courant de l'imprimante et le contexte de production ;
  \item suivre l'encre, le solvant, la viscosité, les températures et les
        heures de fonctionnement ;
  \item estimer le temps restant avant épuisement d'un consommable ;
  \item tenir compte de la date limite d'utilisation et de la durée après
        ouverture ;
  \item enregistrer les défauts et calculer le temps d'arrêt seulement lorsque
        la ligne demande une impression ;
  \item offrir des tableaux de bord distincts pour l'opérateur, la maintenance
        et la direction ;
  \item valider l'ensemble sur un simulateur réaliste avant le raccordement au
        site.
\end{itemize}

\section{Démarche suivie}

La démarche repose sur une séparation claire entre le réseau industriel et la
plateforme de supervision. Litmus Edge constitue le point d'entrée des données.
En aval, une architecture conteneurisée assure le transport MQTT, la validation
des messages, leur stockage temporel et leur présentation dans Grafana.

Le développement a été mené avec un simulateur d'imprimante CIJ afin de disposer
de scénarios reproductibles : production normale, attente sans commande,
consommation des fluides, dérive de viscosité, remplissage et défaut. La même
chaîne pourra ensuite recevoir les messages de Litmus Edge en remplaçant le
mapping d'entrée, sans modifier les tableaux de bord.

\section{Organisation du rapport}

Le chapitre suivant présente Saint-Gobain Sekurit Maroc et le procédé de
fabrication dans lequel s'inscrit le marquage. L'état de l'art conduit ensuite
à la problématique, puis le cahier des charges précise les besoins et les
contraintes. Un chapitre distinct décrit la conduite du projet, ses incréments
et la gestion des risques. La conception, la réalisation et les tableaux de
bord sont présentés avant les essais. Les derniers chapitres traitent de la
sécurité, de l'industrialisation, de l'exploitation et du bilan du travail.
